\documentclass[11pt]{ctexart}

\usepackage{fontspec}
\usepackage{amsmath}
\usepackage{amssymb}
\usepackage{hyperref}
\usepackage[margin=1.2in]{geometry}
\usepackage{booktabs}
\usepackage{enumitem}
\usepackage{parskip}
\usepackage{setspace}
\usepackage{graphicx}
\usepackage{xcolor}
\usepackage{tcolorbox}

\setstretch{1.15}
\setcounter{secnumdepth}{0}
\date{}

\hypersetup{
    colorlinks=true,
    linkcolor=blue,
    urlcolor=blue,
    citecolor=blue,
}

\title{\LARGE\bfseries QK-Norm 可能恶化 Muon 优化器 LLM 训练}
\author{Vuk Rosić \\ \small \raisebox{-0.2\height}{\includegraphics[width=0.02\textwidth]{/Users/vukrosic/AI Science Projects/AI Blog Writing/.agent/skills/md-to-pdf/github.jpg}} \texttt{vukrosic/qk\_norm\_collapse}}

\begin{document}
\maketitle

\begin{tcolorbox}[colback=gray!10, colframe=gray!50, arc=2mm, boxrule=0.5pt]
\textbf{注意：} 这是一篇探索性的博文，记录了单次种子实验的早期观察结果。这些发现是方向性的信号，而非既定结论。我们在每个条件下运行一个种子，仅训练 50M token（约为 1.5B 模型 Chinchilla 最优值的 0.3\%）。所有主张应解读为“我们观察到了 X”，而不是“X 普遍为真”。
\end{tcolorbox}

\section{主要发现一览}

\begin{figure}[h!]
\centering
\includegraphics[width=\textwidth]{./images/2_key_pr_causal.png}
\caption{键秩崩塌——是 $\gamma$ 还是归一化导致的？}
\end{figure}

Transformer 中的每个注意力头使用 \textbf{128 维空间}来表示 token。理想情况下，模型将其表示分布在所有 128 个维度上。实际上，许多维度的方差收缩到接近于零——模型“崩塌”到一个低维子空间。上图中的 y 轴测量了这 128 个维度中有多少个活跃地承载方差（越高 = 使用的维度越多）。\textbf{重要提示：}接近零的方差并不意味着维度无用——它仍然可能通过微小的变化编码关键特征。PR（参与率）衡量的是几何分布，而非信息含量。

\textbf{QK-Norm} 是一种对注意力向量做两件事的技术：(1) 它对向量进行\textbf{归一化}（控制其幅度），以及 (2) 它将每个维度乘以一个\textbf{可学习权重 $\gamma$}——本质上是给模型的每个维度一个音量旋钮。我们想知道：当 QK-Norm 起作用时，实际上是哪一部分在起作用？

如果您需要 QK-Norm 和 $\gamma$ (gamma) 的解释，请先向下滚动。

我们将一个 \textbf{1.5B 参数}（完整规格见下文）的语言模型训练了三次——一次使用完整的 QK-Norm，一次使用归一化但将 $\gamma$ 旋钮锁定在 1.0（因此模型无法调整它们），还有一次完全不使用 QK-Norm。以下是我们的发现：

\begin{enumerate}
\item \textbf{$\gamma$ 旋钮——而非归一化——似乎控制着有多少维度承载方差。} 橙色线（完整 QK-Norm）恢复到 65.7 个维度。紫色线（归一化但锁定 $\gamma$）仅恢复到 59.6 —— \textit{比}青色线（完全无 QK-Norm，63.3）\textit{更差}。在这个单种子实验中，锁定 $\gamma$ 旋钮使模型比不归一化\textit{更糟}。

\item \textbf{但是归一化——而非 $\gamma$——是让模型学习得更好的原因。} 两个归一化模型（橙色和紫色）达到了几乎相同的损失，并且都击败了无 QK-Norm 模型。在这个小规模、未充分训练的模型上，$\gamma$ 旋钮根本不影响学习速度。

\begin{itemize}
\item 即使头内部的某些维度变得接近于零（包含较少数据），模型似乎仍然能很好地学习预测下一个 token。
\end{itemize}

\item \textbf{每个模型都经历了相同的“崩塌-恢复”模式。} 所有三个模型都以 \textasciitilde87 个在用维度开始，暴跌至 \textasciitilde51--55，然后恢复。这种 U 形出现在所有三个条件下，表明这可能是早期 Transformer 训练的一个普遍特征（尽管这需要进一步调查）。

\item \textbf{我们无法将 PR 与模型质量联系起来。} 这是最重要的警告：具有\textit{最少}活跃维度的模型（冻结 $\gamma$，PR=59.6）实现了\textit{最佳}损失 (3.614)。\textbf{在我们的实验中，更高的 PR 并不意味着更好的模型。} $\gamma$ 控制一种几何属性（方差分布的均匀程度），我们尚无法将其与下游性能联系起来。崩塌的维度是代表浪费的计算还是高效的压缩仍然是一个悬而未决的问题——答案很可能是“模型不需要所有 128 个维度”。
\end{enumerate}

下面我们将定义每个技术概念，然后逐步分析图表。

\subsection{您首先需要了解的概念}

\subsubsection{什么是 $\gamma$ (gamma)？}

在 Transformer 注意力头中，每个 token 被投影到一个 128 维的\textbf{键 (Key)}向量。 \textbf{QK-Norm} 对此向量应用 RMSNorm，它做两件事：

\begin{enumerate}
\item \textbf{归一化} - 将每个元素除以整个向量的均方根，从而使得向量具有受控的幅度。
\item \textbf{按 $\gamma$ 缩放} - 将 128 个元素中的每一个乘以其自己的可学习权重 $\gamma_j$。
\end{enumerate}

$$\text{RMSNorm}(\mathbf{x})_j = \gamma_j \cdot \frac{x_j}{\sqrt{\frac{1}{128}\sum_{i=1}^{128} x_i^2}}$$

将 $\gamma$ 想象为一个\textbf{每维度的音量旋钮}。128 个维度中的每一个都有自己的旋钮，初始化为 1.0（全部相等）。随着训练的进行，模型可以调高某些旋钮（放大该维度）并调低其他旋钮（抑制它），从而直接控制哪些维度重要。

\textbf{具体示例}，一个 4 维向量 $\mathbf{x} = [2, 4, 6, 8]$：
\begin{itemize}
\item RMS $= \sqrt{(4 + 16 + 36 + 64)/4} = \sqrt{30} \approx 5.48$
\item 归一化后：$[0.365,\ 0.730,\ 1.095,\ 1.461]$ (所有值重新缩放)
\item 现在假设训练学习到了 $\gamma = [1.5, 1.0, 0.5, 0.1]$:
\begin{enumerate}
\item 维度 1: $0.365 \times 1.5 = 0.548$ $\leftarrow$ \textbf{被放大}
\item 维度 2: $0.730 \times 1.0 = 0.730$ $\leftarrow$ 未改变
\item 维度 3: $1.095 \times 0.5 = 0.548$ $\leftarrow$ \textbf{被抑制}
\item 维度 4: $1.461 \times 0.1 = 0.146$ $\leftarrow$ \textbf{几乎被移除}
\end{enumerate}
\item 最终结果: $[0.55,\ 0.73,\ 0.55,\ 0.15]$ - $\gamma$ 重塑了哪些维度承载信息。
\end{itemize}

\subsubsection{什么是参与率 (PR/Participation Ratio)？}

PR 衡量在键向量中\textbf{这 128 个维度中有多少个承载了显著的方差}。它回答了：“模型是将方差分布在许多维度上，还是将其集中在少数几个维度上？”

\begin{itemize}
\item \textbf{PR = 128} $\rightarrow$ 所有维度贡献均等（最大多样性，满秩）
\item \textbf{PR = 1} $\rightarrow$ 严格来说意味着只有一个奇异值非零，其他全为零。实际上，奇异值永远不会完全为零——它们只是非常小。因此，真实模型可能会在严重崩塌中显示 PR = 2--3，这意味着方差绝大多数集中在 2--3 个维度中，而其余 \textasciitilde125 个维度的方差可忽略不计（但不为零）。
\item \textbf{PR $\approx$ 60} (我们观察到的) $\rightarrow$ 128 个维度中大约有 60 个承载显著方差
\end{itemize}

当 PR 下降时，方差集中到更少的维度——许多键维度的方差收缩到接近零。这被称为\textbf{维度崩塌}（或秩崩塌）。\textbf{重要警告：}低方差维度不一定是无用维度。它仍然可能编码微妙但关键的区别（例如，二进制特征）。PR 衡量的是几何分布，而非信息含量。崩塌的维度是代表浪费的计算还是高效的压缩仍然是一个悬而未决的问题。

\subsubsection{三个实验条件}

我们训练同一个 1.5B 模型三次，仅改变 QK-Norm 的配置方式：

\begin{table}[h!]
\centering
\begin{tabular}{lccp{5cm}}
\toprule
\textbf{条件} & \textbf{归一化？} & \textbf{$\gamma$ 可学习？} & \textbf{目的} \\
\midrule
\textbf{学习 $\gamma$} & \checkmark\ 是 & \checkmark\ 是 & 完整的 QK-Norm —— 真实情况 \\
\textbf{冻结 $\gamma$=1} & \checkmark\ 是 & $\times$\ 否 (锁定在 1.0) & 消融实验 —— 有归一化但\textit{没有}学习缩放 \\
\textbf{无 QK-Norm} & $\times$\ 否 & $\times$\ 否 & 基线 —— 完全没有归一化 \\
\bottomrule
\end{tabular}
\end{table}

\textbf{为什么有三个条件？} 因为 QK-Norm 同时做两件事（归一化 + 按 $\gamma$ 缩放）。仅有两个条件（开/关）我们无法分辨哪部分重要。冻结 $\gamma$=1 条件在保持归一化的同时通过外科手术移除了 $\gamma$ 学习，让我们能够分离原因。

\subsection{图 1：核心结果——是 $\gamma$ 还是归一化驱动了秩？}

\begin{figure}[h!]
\centering
\includegraphics[width=\textwidth]{./images/2_key_pr_causal.png}
\caption{键秩崩塌——是 $\gamma$ 还是归一化导致的？}
\end{figure}

\textbf{此图显示的内容：} 较低的 PR（y 轴）可能表示更多的计算浪费或无用维度，但这需要进一步调查，所以让我们严格定义它：y 轴是\textit{参与率 (PR)}，即承载显著方差的维度数量（满分 128）。越高意味着方差分布越均匀；越低意味着它集中在更少的维度中。x 轴是训练进度（百万 token）。

\textbf{逐步解读：}

\begin{enumerate}
\item \textbf{所有三个模型都从 PR $\approx$ 87 开始}（左上角）。在随机初始化时，128 个维度中的每一个贡献大致相等——模型还没有学到任何东西。

\item \textbf{所有三个模型在 \textasciitilde8M token 时崩塌至 PR $\approx$ 51--55}（U 形曲线底部）。这就是\textit{秩崩塌}——模型在开始学习时破坏了其初始的随机结构。无论是否使用 QK-Norm，这都会发生。

\item \textbf{在崩塌底线之后，三条线分叉了}——这就是实验揭示答案的地方：
\begin{itemize}
\item \textbf{学习 $\gamma$ (橙色)} 恢复得最快且最高，达到 PR = \textbf{65.7}
\item \textbf{无 QK-Norm (青色)} 恢复到 PR = \textbf{63.3} —— 令人惊讶的是，排名第二
\item \textbf{冻结 $\gamma$=1 (紫色)} 恢复最少，仅达到 PR = \textbf{59.6}
\end{itemize}

\item 如果归一化是恢复的关键机制，那么两个归一化模型（学习 $\gamma$ 和冻结 $\gamma$）应该恢复得相似。相反，移除 $\gamma$ 学习（冻结）使其比完全没有归一化\textit{更差}。\textbf{在这个实验中，可学习的 $\gamma$ 参数——而不是归一化——似乎是驱动秩恢复的原因。}（警告：这是单种子结果；6 个单位的差异可能会随不同的种子变窄或变宽。）
\end{enumerate}

\subsection{图 2：损失讲述了一个不同的故事}

\begin{figure}[h!]
\centering
\includegraphics[width=\textwidth]{./images/1_loss.png}
\caption{训练和验证损失}
\end{figure}

\textbf{该图显示的内容：} 相同 50M token 期间的训练和验证损失（越低 = 语言建模越好）。

\textbf{逐步解读：}

\begin{enumerate}
\item \textbf{两个归一化模型（学习 $\gamma$ 和冻结 $\gamma$）实现了几乎相同的损失} —— 训练大约 3.62，验证大约 3.44。它们几乎完美重叠。

\item \textbf{无 QK-Norm 模型的损失明显更高} —— 训练 3.68，验证 3.51。归一化显然有助于模型更好地学习。

\item \textbf{但这正是关键见解所在：} 将此与上面的图 1 进行比较。对于\textit{损失}，无论 $\gamma$ 是学习的还是冻结的，归一化的帮助都是均等的。对于\textit{秩}，只有学习 $\gamma$ 有帮助——冻结 $\gamma$ 实际上有害。\textbf{这表明损失和秩是由不同的机制驱动的：}
\begin{itemize}
\item \textbf{归一化 $\rightarrow$ 稳定梯度 $\rightarrow$ 降低损失} (不需要 $\gamma$)
\item \textbf{学习 $\gamma$ $\rightarrow$ 选择性地放大/抑制维度 $\rightarrow$ 维持维度多样性} (需要 $\gamma$)
\end{itemize}

\item \textbf{出现了一个悬而未决的问题：} 冻结 $\gamma$=1 模型具有\textit{最低}的 PR (59.6) 但实现了\textit{最佳}的损失 (3.614)。如果崩塌的维度包含关键信息，我们会预期更差的损失。这表明 (a) 在这个训练阶段，这些维度确实是多余的，(b) 50M token 太早以至于秩的差异无法在损失中体现出来，或者 (c) PR 捕获的几何属性并不直接映射到任务相关信息。我们尚无法确定哪种解释是正确的。
\end{enumerate}

\subsection{解释三个主要观察结果}

\begin{tcolorbox}[colback=gray!10, colframe=gray!50, arc=2mm, boxrule=0.5pt]
\textbf{观察 1：“$\gamma$ 驱动维度多样性，归一化驱动损失。”} 这两个效应似乎是独立的。无论 $\gamma$ 是学习的还是冻结的，归一化都能同等地降低损失。学习 $\gamma$ 会提高 PR —— 但\textbf{我们尚无法将更高的 PR 与更好的模型质量联系起来}（冻结 $\gamma$ 具有\textit{最低}的 PR 但\textit{最佳}的损失）。
\end{tcolorbox}

\begin{tcolorbox}[colback=gray!10, colframe=gray!50, arc=2mm, boxrule=0.5pt]
\textbf{观察 2：“PR 遵循 U 形轨迹。”} 所有三个条件都显示出相同的崩塌-恢复弧线，这表明这可能是早期 Transformer 训练的一个普遍特征，而不是 QK-Norm 的产物。
\end{tcolorbox}

\begin{tcolorbox}[colback=gray!10, colframe=gray!50, arc=2mm, boxrule=0.5pt]
\textbf{观察 3：“没有 $\gamma$ 的归一化似乎对秩产生反作用。”} 将 $\gamma$ 冻结在 1 产生的 PR 比完全没有 QK-Norm 还要\textit{差}。假设：归一化将键约束在一个球面上；如果没有 $\gamma$ 来选择性地拉伸维度，这种约束会限制谱多样性。这尚未得到验证。
\end{tcolorbox}

\subsection{额外观察：深度-PR 梯度反转}

\begin{tcolorbox}[colback=gray!10, colframe=gray!50, arc=2mm, boxrule=0.5pt]
在 500K token 时，标准梯度成立（浅层 $\approx$ PR 122，深层 $\approx$ 106）。到了 50M token 时，即使情况\textbf{反转} —— 更深的层发展出\textit{更高}的 PR。\textbf{警告：} 这两个测量来自具有略微不同配置的独立实验；严格的分析需要在一个运行中连续跟踪每层。
\end{tcolorbox}

\section{1. 概述}

我们研究了 \textbf{QK-Norm} 如何影响 Transformer 注意力头中的维度崩塌。我们在三种条件（学习 $\gamma$，冻结 $\gamma$=1，无 QK-Norm）下训练了一个 \textbf{1.5B 参数} Gemma 风格的 LLM \textbf{\textasciitilde49M token}，以区分归一化与可学习尺度参数的作用。这是单种子消融实验——发现是初步的。

\subsection{架构与训练配置}

\begin{table}[h!]
\centering
\begin{tabular}{ll}
\toprule
\textbf{参数} & \textbf{值} \\
\midrule
\textbf{模型大小} & \textasciitilde1.58B 参数 \\
$d_\text{model}$ & 2048 \\
层数 & 32 \\
注意力头 & 16 ($d_k = 128$) \\
KV 头 & 8 (GQA) \\
$d_\text{ff}$ & 8192 \\
\textbf{训练预算} & 每个条件 \textbf{\textasciitilde49M Token} \\
\textbf{硬件} & NVIDIA H100 80GB \\
批大小 & 8 \\
梯度累积 & 4 \\
有效批大小 & $8 \times 4 \times 2048 = 65{,}536$ tokens/step \\
优化器 & Muon + AdamW \\
数据集 & FineWeb/Cosmopedia Mix (1B subset) \\
\bottomrule
\end{tabular}
\end{table}

\section{2. 结果摘要}

\begin{table}[h!]
\centering
\begin{tabular}{lccc}
\toprule
\textbf{指标} & \textbf{学习 $\gamma$} & \textbf{冻结 $\gamma$=1} & \textbf{无 QK-Norm} \\
\midrule
\textbf{初始平均 PR} & 86.85 & 86.85 & 86.42 \\
\textbf{最小平均 PR} & 55.17 (7.9M) & 51.84 (9.8M) & 51.28 (7.9M) \\
\textbf{最终平均 PR (49M)} & \textbf{65.70} & 59.55 & 63.25 \\
\textbf{最终训练损失} & \textbf{3.618} & \textbf{3.614} & 3.683 \\
\textbf{最终验证损失} & \textbf{3.440} & \textbf{3.439} & 3.505 \\
\bottomrule
\end{tabular}
\end{table}

\section{4. 数学背景}

\subsection{4.1 QK-Norm}

在标准注意力中，token 被投影为查询和键：$Q = XW_Q,\ K = XW_K$。使用 QK-Norm，在 RoPE 之前对 Q 和 K 应用 RMSNorm：

$$Q = \text{RoPE}\!\left(\text{RMSNorm}(XW_Q)\right), \quad K = \text{RoPE}\!\left(\text{RMSNorm}(XW_K)\right)$$

$$\text{RMSNorm}(\mathbf{x})_j = \frac{\gamma_j \cdot x_j}{\text{RMS}(\mathbf{x})}, \quad \text{RMS}(\mathbf{x}) = \sqrt{\frac{1}{d_k}\sum_{i=1}^{d_k} x_i^2}$$

每个 RMSNorm 都有其自己的 $\gamma \in \mathbb{R}^{128}$，初始化为 1。共有 $32 \times 2 = 64$ 个这样的向量（总共 8,192 个额外参数）。随着训练的进行，个别 $\gamma_j$ 值发散——放大某些维度并抑制其他维度。

RoPE 在 RMSNorm \textbf{之后}应用，因此 $\gamma$ 作用于位置前的表示。

\subsection{4.2 为什么这对秩很重要}

PR 通过 SVD 奇异值 $\sigma_1 \geq \cdots \geq \sigma_{128}$ 衡量方差在 128 个键维度上的分布均匀程度：

$$\text{PR}(K) = \frac{\left(\sum_{i=1}^{d_k} \sigma_i\right)^2}{\sum_{i=1}^{d_k} \sigma_i^2}$$

PR = 128 意味着所有维度贡献均等；PR = 1 意味着完全崩塌。每个条件对秩的影响不同：没有 QK-Norm，秩仅通过 $W_K$ 控制；有学习 $\gamma$，模型具有直接的每维度旋钮；有冻结 $\gamma$=1，归一化约束了几何形状，但没有每维度的灵活性。

\section{5. $\gamma$ 动力学}

$\gamma$ 变异系数 (CV) 揭示了模型如何使用可学习参数。高 CV 层（L0: 0.169, L15: 0.155, L27: 0.172）显示出积极的维度分化，并与更高的 PR 恢复相关。低 CV 层（L5--L10: 0.07--0.08）显示出最小的分化。

\subsection{5.1 最终 $\gamma$ 值：模型学到了什么？}

\textbf{$\gamma$ 值大多 $>$ 1.0 —— 模型放大而不是抑制。} 全局平均 $\gamma \approx 1.19$。只有层 0 的平均 $\gamma < 1.0$ (0.913)。

\begin{table}[h!]
\centering
\begin{tabular}{lccp{5cm}}
\toprule
\textbf{层区域} & \textbf{层} & \textbf{平均 $\gamma$} & \textbf{解释} \\
\midrule
\textbf{浅层} & 0--7 & 1.084 & 最接近初始化，层 0 是离群点 \\
\textbf{中层} & 8--23 & 1.298 & 最强的放大 \\
\textbf{深层} & 24--31 & 1.238 & 略少于中层 \\
\bottomrule
\end{tabular}
\end{table}

\textbf{层 0 是独特的不对称的。} 前 64 个维度（高频 RoPE）的平均 $\gamma = 1.017$，而后 64 个维度（低频 RoPE）被抑制到 0.810。所有 14 个 $\gamma < 0.7$ 的维度都落在后 64 个中。这表明层 0 学习去强调低频位置信息。

\textbf{$\gamma$ 分化与 PR 恢复相关。} 具有最高 CV 的层往往显示出最独特的 PR 行为；$\gamma$ 几乎一致的层显示出最小的维度分化。

\textbf{关键局限性：} 我们只有最终检查点的 $\gamma$。而在训练过程中跟踪 $\gamma$ 轨迹将揭示维度选择发生的时间——这是一个明确的下一步。

\section{6. 局限性}

\begin{enumerate}
\item \textbf{单种子}：效应大小（2--6 个 PR 单位）可能会随不同的种子发生变化或反转。
\item \textbf{早期训练}：50M token 约为 Chinchilla 最优值的 0.3%；我们仅观察到早期阶段的动力学。
\item \textbf{仅键分析}：我们探测了键但没有探测查询；$QK^T$ 的有效秩可能讲述一个不同的故事。
\item \textbf{Muon 优化器}：发现可能是 Muon 特定的产物。
\item \textbf{PR $\neq$ 质量}：冻结 $\gamma$ 具有\textit{最低}的 PR 但\textit{最佳}的损失。在我们的数据中，更高的 PR 并不意味着更好的模型。
\item \textbf{未保存奇异值分布}：我们无法区分逐渐变细和急剧截断。
\item \textbf{无下游评估}：我们仅测量损失，没有基准测试。
\end{enumerate}

\subsection{未来工作}

\begin{itemize}
\item \textbf{多种子} ($\geq$3) 用于方差估计——最重要的下一步。
\item \textbf{将训练扩展到 250M--1B+ token} 以测试平台期稳定性以及 PR/损失分离是否解决。
\item \textbf{测量查询 PR} 并直接测量 $QK^T$ 的有效秩。
\item \textbf{保存完整的奇异值分布} 并在训练过程中\textbf{跟踪单独的 $\gamma$ 轨迹}。
\item \textbf{纯 AdamW 基线} 以隔离 Muon 特定效应。
\item \textbf{下游基准测试评估}。
\end{itemize}

\end{document}
